% Seccion 1 de la rubrica: introduccion y justificacion. Abre con \section y
% aloja como \subsection a revi_medio_ambiente.tex, res_sima.tex y
% contexto_estaciones.tex, que se \input inmediatamente despues en main.tex.
% El objeto de investigacion (el ozono) vivia en objetivo.tex; se fusiono aqui
% porque es justificacion, no objetivo.

\section{Introducción y justificación}
\label{sec:introduccion}

La calidad del aire en Nuevo León es un factor determinante de la salud
pública: la exposición diaria a concentraciones elevadas de contaminantes
representa un riesgo acumulativo para la población. Las mediciones continuas
del Sistema Integral de Monitoreo Ambiental (SIMA), descrito en la
\cref{sec:sima}, permiten caracterizar la atmósfera del Área Metropolitana
de Monterrey (AMM) y son la fuente de este trabajo. Para reportar la calidad
del aire, SIMA convierte las mediciones horarias en un índice categórico con
un color asociado a su valor máximo \parencite{aireNL2026}.

De los contaminantes que mide la red, el ozono troposférico es el que
justifica este estudio: es secundario, no se emite sino que se forma, y su
formación depende a la vez de los precursores químicos y de las condiciones
meteorológicas. Esa doble dependencia es lo que vuelve pertinente un análisis
multivariado: ninguna variable por sí sola explica un día de excedencia.

\subsection{El ozono troposférico}
\label{sec:objetivo}

El ozono \ce{O3} es una sustancia compuesta por tres átomos de oxígeno que se
encuentra tanto en la atmósfera superior (ozono estratosférico) como a nivel
del suelo (ozono troposférico). El estratosférico se produce de manera
natural por la interacción entre el oxígeno molecular \ce{O2} y los rayos
ultravioleta, y atenúa la exposición de la superficie a esa radiación
\parencite{Bernhard2023-lk}. El troposférico, en cambio, es el principal
ingrediente del \textit{smog} y se forma por transformaciones químicas entre
los óxidos de nitrógeno \ce{NOx} y los compuestos orgánicos volátiles (COV)
con el calor y la luz solar.

Áreas metropolitanas de Nuevo León con elevada actividad industrial han
presentado niveles de \ce{O3} y \ce{PM10} que superan los límites de la
NOM-020-SSA1-2021 \parencite{nom020ssa1} y la NOM-025-SSA1-2014
\parencite{nom025ssa1}. Estas tendencias se documentan desde hace más de
30 años, con herramientas que miden tanto el ozono (fotometría ultravioleta)
como sus precursores \ce{NO} y \ce{NO2} (agrupados como \ce{NO_x}, por
quimioluminiscencia), cuya relación fotoquímica es clave para entender la
formación del contaminante secundario \parencite{SanfordSillman}.

El AMM no es la excepción: entre 1993 y 2014 se registró un aumento
sostenido de las concentraciones de ozono troposférico, en contraste con la
tendencia descendente de la Ciudad de México en el mismo periodo, atribuido
al crecimiento industrial y vehicular sin un control equivalente sobre los
precursores y a un papel meteorológico modulado por el cambio climático
\parencite{atmos15070779}. El comportamiento se agrava por la topografía del
valle donde se asienta el AMM, rodeado de sierras que favorecen el
estancamiento de masas de aire durante episodios de alta presión, inversión
térmica y viento débil, factores que reducen la dispersión y prolongan la
exposición de la población.
